\chapter{Protocoles I~: FIX}\label{ch:ll:protocols-i-fix}

Un ordre envoyé ce matin par un fonds à son courtier a très probablement voyagé sous la forme d'une ligne de texte~: \texttt{35=D} pour un nouvel
ordre, \texttt{55=ESZ6} pour l'instrument, \texttt{54=1} pour un achat, \texttt{38=5} pour la quantité, chaque couple terminé par un
caractère de contrôle. Le format a été rédigé pour la première fois en 1992 afin que Fidelity Investments et Salomon Brothers puissent échanger
des informations de négociation d'actions par voie électronique, et il est devenu le langage commun des ordres entre le buy side, les courtiers et
de nombreuses plateformes de négociation~: le protocole Financial Information eXchange. Il est ancien, verbeux, lent à analyser naïvement, et omniprésent. Ce
chapitre construit un \omterm{def:ll:protocols-i-fix:engine}{moteur FIX}~: le codec qui lit et écrit le texte sans le recopier, la couche session qui garde
identiques les vues de deux sociétés sur une conversation malgré les déconnexions, et les mesures qui disent ce que coûte le protocole.

\section{Encodage étiquette-valeur et le message}

\begin{definition}[protocole FIX, encodage étiquette-valeur]\label{def:ll:protocols-i-fix:fix}
Le \emph{protocole FIX}\index{protocole FIX} est une norme pour les messages relatifs au trading (ordres, exécutions, cotations, données
de marché, allocations) échangés entre sociétés, maintenue par le FIX Trading Community, avec une couche session qui les délivre dans l'ordre
et sans perte. Dans son \emph{encodage étiquette-valeur}\index{encodage étiquette-valeur} (tag-value), un message est une suite de champs
\texttt{tag=value}, chacun terminé par l'octet SOH (valeur 1), où l'étiquette est un nombre tiré du dictionnaire du protocole et la
valeur est du texte.
\end{definition}

Trois champs ouvrent chaque message, dans un ordre fixe~: BeginString (étiquette 8, la version du protocole), BodyLength (9) et MsgType (35).
Un seul le referme~: CheckSum (10), dont le champ et le SOH final marquent la fin du message. Entre eux viennent le reste
de l'en-tête standard (identifiants de l'émetteur et du destinataire 49 et 56, numéro de séquence 34, heure d'envoi 52 en UTC) et le
corps. BodyLength compte les octets depuis le début du champ MsgType jusqu'au champ CheckSum non compris, délimiteurs
inclus~; CheckSum est la somme de tous les octets qui le précèdent modulo 256, écrite sur trois chiffres.
La \Cref{fig:ll:protocols-i-fix:anatomy} détaille l'exemple que la description habituelle du protocole développe, un logon
dont le corps mesure 65 octets et dont la somme de contrôle vaut 062~; les tests de la construction reproduisent les deux nombres.

\begin{omfigure}
\resizebox{\linewidth}{!}{%
\begin{tikzpicture}[font=\scriptsize,f/.style={draw=gray!60,inner sep=2pt,minimum height=0.5cm,font=\ttfamily\scriptsize}]
\node[f,fill=gray!10,anchor=west] (a) at (0,0) {8=FIX.4.2|};
\node[f,fill=gray!10,anchor=west] (b) at (a.east) {9=65|};
\node[f,fill=omDef!20,anchor=west] (c) at (b.east) {35=A|};
\node[f,fill=omDef!12,anchor=west] (d) at (c.east) {49=SERVER|};
\node[f,fill=omDef!12,anchor=west] (e) at (d.east) {56=CLIENT|};
\node[f,fill=omDef!12,anchor=west] (g) at (e.east) {34=177|};
\node[f,fill=omDef!12,anchor=west] (h) at (g.east) {52=20090107-18:15:16|};
\node[f,fill=omThm!12,anchor=west] (i) at (h.east) {98=0|};
\node[f,fill=omThm!12,anchor=west] (j) at (i.east) {108=30|};
\node[f,fill=gray!10,anchor=west] (k) at (j.east) {10=062|};
\draw[decorate,decoration={brace,amplitude=4pt},thick] (c.north west) ++(0,0.08) -- ([yshift=0.08cm]j.north east)
  node[midway,above=5pt]{BodyLength~: 5 + 10 + 10 + 7 + 21 + 5 + 7 = 65 octets, de \texttt{35=} jusqu'à \texttt{10=}};
\draw[decorate,decoration={brace,amplitude=4pt,mirror},thick] ([yshift=-0.08cm]a.south west) -- ([yshift=-0.08cm]j.south east)
  node[midway,below=5pt]{CheckSum~: chaque octet avant \texttt{10=}, délimiteurs inclus, totalise 4\,158~; $4\,158 \bmod 256 = 62$};
\node[anchor=north west,align=left,font=\footnotesize] at (0,-1.35) {gris~: encadrement (les deux premiers champs, dernier champ)~; bleu~:
  MsgType et le reste de l'en-tête standard~; rouge~: le corps d'un logon (méthode de chiffrement 98, intervalle de heartbeat 108)};
\end{tikzpicture}}

\omcaption{Un logon FIX~4.2 sous forme de chaîne d'octets, où \texttt{|} désigne le délimiteur SOH. BodyLength et CheckSum sont
calculés sur les plages représentées~; un récepteur les recalcule et rejette le message si l'un ou l'autre diffère.}\label{fig:ll:protocols-i-fix:anatomy}
\end{omfigure}

La somme de contrôle est un contrôle faible, et savoir à quel point il est faible dit au lecteur à quoi il sert.

\begin{proposition}[Ce que détecte la somme de contrôle FIX]\label{prop:ll:protocols-i-fix:checksum}
Modifier un seul octet d'un message change toujours sa somme de contrôle. Échanger deux octets, ou deux champs, ne la change jamais~;
pas davantage tout ensemble de modifications dont les valeurs ont pour somme un multiple de 256.
\end{proposition}

\begin{proof}
Remplacer l'octet $b$ par $b' \ne b$ modifie la somme de $b' - b$, avec $0 < |b' - b| < 256$, qui n'est pas un multiple de 256.
Une somme ne dépend pas de l'ordre de ses termes, donc une permutation la laisse inchangée~; et des modifications $\delta_1, \dots, \delta_k$
laissent la somme modulo 256 inchangée exactement lorsque $\sum_j \delta_j \equiv 0 \pmod{256}$.
\end{proof}

Le transport sous-jacent, TCP, porte déjà ses propres sommes de contrôle sur chaque segment. La somme de contrôle FIX et la longueur du corps sont
donc moins une protection contre les altérations sur le réseau qu'un contrôle de l'encadrement~: un programme qui a perdu sa place dans un
flux d'octets, ou un message tronqué ou recollé par une passerelle défaillante, est détecté au message suivant. L'encadrement lui-même
repose sur les deux premiers champs~: une fois \texttt{8=} et \texttt{9=} lus, la longueur totale du message est connue, le corps plus les
sept octets de \texttt{10=nnn} et de son délimiteur.

\section{La couche session~: numéros de séquence, heartbeats, demandes de renvoi et comblement de trou}

\begin{definition}[session FIX, heartbeat, demande de renvoi, comblement de trou]\label{def:ll:protocols-i-fix:session}
Une \emph{session FIX}\index{session FIX} est la conversation bidirectionnelle entre deux parties identifiées par leurs identifiants d'émetteur et
de destinataire, du logon au logout, au cours de laquelle chaque côté numérote les messages qu'il envoie 1, 2, 3, \dots\ et les conserve
afin de pouvoir les renvoyer. Un \emph{heartbeat}\index{heartbeat} est un message envoyé lorsqu'un côté n'a rien émis pendant
un intervalle convenu, afin que le silence puisse être distingué d'une connexion morte. Une \emph{demande de renvoi}\index{demande de renvoi} (resend request) demande à
l'autre côté de renvoyer une plage de ses messages, après qu'un trou dans leurs numéros de séquence a été détecté. Un
\emph{comblement de trou}\index{comblement de trou} (gap fill) est un message de réinitialisation de séquence, envoyé en réponse à une demande de renvoi, qui remplace une suite de messages
administratifs qui ne seront pas renvoyés, en faisant avancer au-delà d'eux le numéro de séquence attendu par le récepteur.
\end{definition}

Les numéros de séquence (le livre~1, chapitre~28, pour les données de marché) constituent à eux seuls le mécanisme de la couche session. Chaque côté tient deux
compteurs~: le numéro de son prochain message sortant et celui qu'il attend ensuite de l'autre côté. Un message portant le
numéro attendu est traité et le compteur avance. Un numéro supérieur signifie que des messages ont été perdus~: le récepteur envoie une
ResendRequest (MsgType 2) à partir du numéro attendu jusqu'à «~l'infini~» (EndSeqNo 0), et retient les messages suivants jusqu'à ce que le trou soit
comblé. Un numéro inférieur est soit un doublon, marqué par PossDupFlag (43) et ignoré, soit une erreur fatale~: les deux côtés ne
s'accordent plus, et le récepteur se déconnecte. La réponse à une \omterm{def:ll:protocols-i-fix:session}{demande de renvoi} contient de nouveau les messages applicatifs, marqués
PossDupFlag avec leur heure d'envoi d'origine (122), et remplace les messages administratifs (\omterm{def:ll:protocols-i-fix:session}{heartbeats}, demandes de test, \omterm{def:ll:protocols-i-fix:session}{demandes de renvoi}
elles-mêmes) par des messages SequenceReset avec GapFillFlag (123) positionné~: un \omterm{def:ll:protocols-i-fix:session}{heartbeat} datant d'une demi-heure ne signifie plus rien
maintenant. La \Cref{fig:ll:protocols-i-fix:sequence} dessine la conversation de référence de la construction autour de son trou, et
le \cref{lst:ll:protocols-i-fix:trace} reprend la même conversation telle que le moteur la trace.

\begin{omfigure}
\begin{tikzpicture}[font=\scriptsize,>=Stealth]
\def\L{0}\def\R{7.5}
\node[font=\footnotesize\bfseries] at (\L,0.45) {SOCIÉTÉ (nous)};
\node[font=\footnotesize\bfseries] at (\R,0.45) {COURTIER};
\draw[gray,thick] (\L,0.15) -- (\L,-9.6);
\draw[gray,thick] (\R,0.15) -- (\R,-9.6);
\newcommand{\msgR}[3]{\draw[->,#3] (\L,#1) -- node[above,sloped]{#2} (\R,#1-0.25);}
\newcommand{\msgL}[3]{\draw[<-,#3] (\L,#1-0.25) -- node[above,sloped]{#2} (\R,#1);}
\msgR{-0.1}{Logon, seq 1}{omDef}
\msgL{-0.75}{Logon, seq 1}{omDef}
\msgR{-1.4}{NewOrderSingle, seq 2}{omDef}
\msgL{-2.05}{ExecutionReport (new), seq 2}{omDef}
\msgL{-2.7}{Heartbeat, seq 3}{gray}
\draw[omThm,dashed,->] (\R,-3.3) -- (5.6,-3.45) node[left,font=\scriptsize]{perdus~: seq 4 (exécution partielle), seq 5 (heartbeat)};
\msgL{-4.0}{ExecutionReport (filled), seq 6}{omDef}
\node[anchor=east,omThm,align=right] at (-0.15,-4.25) {trou~: attendu 4,\\ retient séq. 6};
\msgR{-4.65}{ResendRequest 7=4 16=0, seq 3}{omThm,thick}
\msgL{-5.3}{ExecutionReport (partial), seq 4, PossDup}{omDef}
\msgL{-5.95}{SequenceReset-GapFill, seq 5, NewSeqNo 6}{omMeth}
\node[anchor=east,omMeth,align=right] at (-0.15,-6.2) {livrer 4, puis le\\ 6 retenu~; attendre 7};
\msgL{-6.6}{ExecutionReport, seq 6 again, PossDup: ignored}{gray}
\msgR{-7.25}{OrderCancelRequest, seq 4}{omDef}
\msgL{-7.9}{ResendRequest 7=3 16=0 (the broker lost ours)}{omThm}
\msgR{-8.55}{GapFill seq 3 (our resend request), NewSeqNo 4}{omMeth}
\msgR{-9.2}{OrderCancelRequest again, seq 4, PossDup}{omDef}
\end{tikzpicture}

\omcaption{La conversation de référence de la construction autour de son trou. Deux messages du courtier sont perdus~; le message suivant révèle le
trou, que comble une \omterm{def:ll:protocols-i-fix:session}{demande de renvoi} (resend request)~: le message applicatif revient marqué comme doublon possible, le \omterm{def:ll:protocols-i-fix:session}{heartbeat} est remplacé par un \omterm{def:ll:protocols-i-fix:session}{comblement de trou} (gap fill), et le message retenu est livré
dans l'ordre. Plus tard, le courtier réclame les nôtres, et notre message administratif est à son tour comblé par un \omterm{def:ll:protocols-i-fix:session}{comblement de trou}.
est à son tour comblé.}\label{fig:ll:protocols-i-fix:sequence}
\end{omfigure}

\omcode{code/low-latency/15-protocols-i-fix/python/trace_short.txt}{7}{22}{La conversation de référence autour du trou, telle que
\texttt{ll\_fix.short} affiche la trace du moteur (la trace complète est \texttt{expected\_trace.txt} dans les données de la construction).}\label{lst:ll:protocols-i-fix:trace}

Les \omterm{def:ll:protocols-i-fix:session}{heartbeats} rendent le silence significatif. À l'ouverture de session, les deux parties conviennent d'un intervalle, HeartBtInt (108), ici de 30 secondes. Une partie qui n'a rien
envoyé pendant cette durée émet un \omterm{def:ll:protocols-i-fix:session}{Heartbeat} (0). Une partie qui n'a rien reçu pendant un peu plus longtemps (la construction attend 1.2
intervalle) émet un TestRequest (1) portant un identifiant que le \omterm{def:ll:protocols-i-fix:session}{heartbeat} de réponse doit reprendre~; si cette réponse n'arrive pas dans un
intervalle supplémentaire, elle déclare la connexion morte. Le délai de détection est le prix d'un long intervalle~: avec 30 secondes, une connexion
silencieusement morte est repérée après environ 66, pendant lesquelles les messages de la contrepartie ne vont nulle part et doivent être
renvoyés.

\omcode{code/firm/fixengine/cpp/firm_fixengine.hpp}{267}{289}{Le contrôle de séquence du moteur C++~: un trou demande un renvoi et retient le
message~; un doublon est ignoré~; un numéro trop bas met fin à la session.}\label{lst:ll:protocols-i-fix:seq}

\section{La couche applicative~: ordres et rapports d'exécution}

Au-dessus de la session, chaque message applicatif porte un événement métier. Un NewOrderSingle (35=D) porte l'identifiant d'ordre de la société
ClOrdID (11), l'instrument (55), le sens (54), la quantité (38), le type d'ordre (40) et, pour un ordre à cours limité, le prix (44)~; un OrderCancelRequest (35=F)
désigne l'ordre à annuler par son identifiant d'origine (41) et donne à la demande son propre nouvel identifiant. La réponse à chacun est un rapport d'exécution (35=8),
que la spécification utilise «~pour confirmer la réception d'un ordre~», pour confirmer les modifications, pour relayer le statut et les exécutions et pour rejeter des ordres~;
le livre~10, chapitre~26, le décrit du côté de la plateforme. Son OrdStatus (39) parcourt le cycle de vie de l'ordre du livre~7, chapitre~17~:
\texttt{0} nouvel, \texttt{1} partiellement exécuté, \texttt{2} exécuté, \texttt{4} annulé, \texttt{8} rejeté~; LastQty (32) et LastPx (31)
décrivent la dernière exécution, CumQty (14) et LeavesQty (151) les totaux.
Une exécution partielle de 2 contrats sur 5 à 5723.25 se lit donc
\texttt{39=1}, \texttt{32=2}, \texttt{31=5723.25}, \texttt{14=2}, \texttt{151=3}.

La couche session et la couche applicative se rencontrent en un point~: les messages renvoyés. Un rapport d'exécution qui arrive avec PossDupFlag
peut avoir été traité ou non auparavant (l'original a pu arriver et seul son successeur avoir été perdu). L'application doit donc rendre le traitement idempotent,
en reconnaissant une exécution déjà enregistrée à son identifiant d'exécution (17), faute de quoi elle enregistre deux fois la même exécution.
La couche session livre chaque numéro de séquence une seule fois~; elle ne peut pas savoir ce que l'application
a fait de l'original.

\section{Analyser vite}

\begin{definition}[moteur FIX]\label{def:ll:protocols-i-fix:engine}
Un \emph{moteur FIX}\index{moteur FIX} est le logiciel qui implémente la couche \omterm{def:ll:protocols-i-fix:session}{session FIX} pour une application~: il encadre et valide les messages entrants,
tient à jour les numéros de séquence, les \omterm{def:ll:protocols-i-fix:session}{heartbeats} et le stock des messages émis, répond aux \omterm{def:ll:protocols-i-fix:session}{demandes de renvoi}
et remet les messages applicatifs au programme dans l'ordre.
\end{definition}

Un analyseur naïf découpe le message sur le SOH, convertit chaque étiquette en entier et copie chaque valeur dans une table associant étiquette et chaîne~: plusieurs allocations
par champ (chapitre~6), une insertion dans un arbre par champ et une passe par octet. L'analyseur de la construction ne fait rien de tout cela. Il repère chaque délimiteur du message
avec le balayage vectoriel du chapitre~14, analyse chaque étiquette sur place, et ne consigne pour chaque champ que son étiquette, son décalage et sa longueur dans un tableau de taille fixe~:
les valeurs restent là où elles sont, dans le tampon de réception, et sont lues comme des vues de chaînes.
La longueur du corps et la somme de contrôle sont vérifiées dans la même passe.

\omcode{code/firm/fixengine/cpp/firm_fixengine.hpp}{40}{73}{L'analyse sans copie~: les délimiteurs issus du balayage vectoriel, les étiquettes analysées sur place, puis
l'ordre des champs d'encadrement, la longueur du corps et la somme de contrôle.}\label{lst:ll:protocols-i-fix:parse}

La \Cref{fig:ll:protocols-i-fix:parse} mesure les deux sur des rapports d'exécution d'environ 205 octets. La vue sans copie analyse l'un en
environ \qty{80}{\nano\second}, soit une douzaine de millions par seconde sur un cœur~; l'analyseur à base de table prend environ \qty{900}{\nano\second},
onze fois plus~; la référence Python environ \qty{5}{\micro\second}. La construction est symétrique~: le constructeur à tampon fixe, qui écrit lui-même les chiffres
et calcule la longueur et la somme de contrôle sur place, produit un nouvel ordre en environ \qty{60}{\nano\second}, contre environ
\qty{330}{\nano\second} pour le même message assemblé par concaténation de chaînes.

\begin{omfigure}
\begin{tikzpicture}
\begin{axis}[width=11.5cm,height=5.6cm,ybar,bar width=14pt,ymode=log,log origin y=infty,ymin=20,ymax=20000,
  symbolic x coords={view,map,python,builder,concat},xtick=data,
  xticklabels={{analyse~:\\ vue sans\\ copie},{analyse~:\\ std::map},{analyse~:\\ référence\\ Python},{construction~:\\ tampon fixe},{construction~:\\ concaté-\\nation}},
  xticklabel style={align=center,font=\scriptsize},enlarge x limits=0.12,
  ylabel={ns par message (échelle log)},tick label style={font=\small},label style={font=\small},grid=major,
  grid style={gray!25},nodes near coords,every node near coord/.append style={font=\scriptsize},
  point meta=explicit symbolic,table/col sep=comma]
\addplot[fill=omDef!60,draw=omDef] table[x=key,y=ns_per_msg,meta=ns_per_msg]{figdata/low-latency/15-protocols-i-fix/measured_fix.csv};
\end{axis}
\end{tikzpicture}

\omcaption{Analyse et construction de rapports d'exécution FIX d'environ 205 octets (1\,000 messages distincts, médiane de 40 passes).
Mesuré sur un portable (Intel Core Ultra 7 155H) sous WSL2, sans cœurs isolés~; C++ avec GCC~11 en \texttt{-O2}, Python
3.10. Données~: \texttt{bench\_fix.py}.}\label{fig:ll:protocols-i-fix:parse}
\end{omfigure}

L'analyse ne représente que rarement tout le coût d'une connexion FIX~: la couche session écrit chaque message émis dans son stock, souvent sur disque,
avant ou après l'envoi, et le moteur de la contrepartie, le réseau et l'appariement prennent bien plus longtemps que \qty{80}{\nano\second}.
Mais sur une passerelle qui porte tout le flux d'un courtier, ou sur le chemin d'ordres d'un teneur de marché passant par une plateforme qui ne parle que FIX,
l'écart entre \qty{80}{\nano\second} et une microseconde par message sépare un composant qui disparaît du \omterm{def:ll:where-latency-comes-from:budget}{budget de latence} du chapitre~1
d'un composant qui en occupe une part.

\section{Où FIX reste utilisé}

FIX est le langage du routage d'ordres entre sociétés~: des systèmes de gestion des ordres (OMS) des sociétés de gestion vers les courtiers,
des algorithmes des courtiers vers les plateformes, et pour l'accès direct au marché (DMA, livre~1, chapitre~4), où les ordres d'un client
atteignent la plateforme par l'infrastructure du courtier. Certaines plateformes l'acceptent pour la saisie d'ordres et les drop copy, comme le montre
l'encadré daté ci-dessous pour l'une d'elles. Là où la vitesse est le produit, les plateformes sont passées à des protocoles binaires aux champs à décalage fixe (chapitre~16),
et le FIX Trading Community a lui-même normalisé des alternatives binaires~: Simple Binary Encoding pour les messages, et FIXP, un «~protocole point à point léger~»
pour la session, dont l'introduction cite pour base la couche \omterm{def:ll:protocols-i-fix:session}{session FIX}
et des protocoles d'échange comme SoupBinTCP. Le moteur à étiquette-valeur reste
celui dont chaque société de trading a besoin quelque part.

\begin{dated}{2026-09}{La saisie d'ordres FIX aujourd'hui}\label{dat:ll:protocols-i-fix:today}
Coinbase Exchange documente trois façons de négocier (consulté en septembre 2026)~: une interface REST «~pour un trading à plus faible fréquence~»,
une interface de saisie d'ordres FIX «~pour un trading à plus haute fréquence~», et les données de marché en FIX et WebSocket. Sa référence FIX
est FIX 5.0 SP2~; les sessions sont fermées chaque samedi~; et «~resend requests are not supported~»~: chaque connexion ouvre une nouvelle session
avec de nouveaux numéros de séquence, si bien qu'un client qui perd sa connexion retrouve l'état de ses ordres par d'autres moyens,
comme les sessions drop copy de la plateforme. La norme de session FIXP a atteint la version 1.0 en avril 2021.
\end{dated}

\section{Tutoriel~: un moteur FIX en trois langages}\label{tut:ll:protocols-i-fix:tut}

\begin{tutorial}
\textbf{Objectif.} Analyser et construire des messages FIX sans les copier, rejouer une conversation avec un trou à travers la couche session
en Python et en C++, et mesurer les analyseurs. \textbf{État final~:} la \cref{fig:ll:protocols-i-fix:parse}, la trace du
\cref{lst:ll:protocols-i-fix:trace}, et des tests au vert en Python, C++ et Rust.
\begin{tutsteps}
\item \textbf{Vérifier l'exemple publié.} Les tests décodent le logon de la \cref{fig:ll:protocols-i-fix:anatomy} et trouvent un
corps de 65 octets et une somme de contrôle de 062, puis modifient un octet (erreur de somme de contrôle) et la longueur déclarée (erreur de longueur de corps).
\item \textbf{Rejouer la conversation de référence.} \texttt{make\_fix\_fixtures.py} scénarise le côté du courtier (une exécution perdue et un
\omterm{def:ll:protocols-i-fix:session}{heartbeat} perdu, une \omterm{def:ll:protocols-i-fix:session}{demande de renvoi} du courtier, un courtier silencieux)~; \texttt{firm\_fixengine.replay} fait passer notre session dans cette conversation
et écrit la trace~; le test C++ fait passer la session C++ par la même conversation et compare sa trace à celle de Python,
ligne à ligne~; les tests Rust analysent et reconstruisent chaque message, octet par octet.
\item \textbf{Raccourcir la trace} avec \texttt{python ll\_fix.py}, qui affiche le \cref{lst:ll:protocols-i-fix:trace}.
\item \textbf{Mesurer} avec \texttt{python bench\_fix.py}~: les analyseurs et constructeurs C++ sur 1\,000 rapports d'exécution,
et la référence Python sur les mêmes messages.
\end{tutsteps}
\textbf{Pour aller plus loin.} Supprimer le \omterm{def:ll:protocols-i-fix:session}{comblement de trou} de la réponse du courtier dans la conversation et rejouer celle-ci~: la session
attend le message 5 indéfiniment. Faire ensuite répondre la session à un message dont l'étiquette n'est pas un nombre par un Reject (35=3).
\end{tutorial}

\section{Construction~: le moteur FIX}\label{bld:ll:protocols-i-fix:fixengine}

\begin{build}
\bfield{Rôle} Les connexions FIX de la société~: saisie d'ordres auprès des courtiers et des plateformes où FIX est le protocole, drop copies, et les
contreparties de test du chapitre~25. La passerelle d'ordres du chapitre~21 utilise son codec pour les plateformes FIX.
\bfield{Interface} Référence Python \texttt{firm\_fixengine}~: \texttt{checksum}, \texttt{encode}, \texttt{decode},
\texttt{frame}, \texttt{Session} (\texttt{logon}, \texttt{send}, \texttt{on\_bytes}, \texttt{on\_timer}, \texttt{logout},
\texttt{trace}), \texttt{replay}. C++20 \texttt{firm::fix}~: \texttt{View::parse}, \texttt{get}, \texttt{frame},
\texttt{Builder} (\texttt{start}, \texttt{add}, \texttt{finish}), \texttt{Session} avec les mêmes méthodes et la même trace. Rust
\texttt{firm\_fixengine}~: \texttt{View::parse}, \texttt{get}, \texttt{frame}, \texttt{encode}, \texttt{checksum}.
\bfield{Règles} L'analyseur ne copie rien et rejette tout message dont les champs de cadrage, la longueur du corps ou la somme de contrôle sont incorrects~; le
constructeur n'alloue rien~; la session numérote chaque message, conserve chaque message qu'elle envoie, répond aux \omterm{def:ll:protocols-i-fix:session}{demandes de renvoi} par
des messages applicatifs marqués PossDupFlag et par des \omterm{def:ll:protocols-i-fix:session}{comblements de trou} pour les messages d'administration, retient les messages reçus au-delà d'un trou, et
met fin à la session sur un numéro de séquence trop bas sans PossDupFlag.
\bfield{Tests d'acceptation} \texttt{code/firm/fixengine/}~: la longueur du corps et la somme de contrôle de l'exemple publié~; les corruptions de la somme de contrôle,
de la longueur du corps et de l'ordre des champs détectées~; le cadrage de flux concaténés et partiels~; la trace de la conversation de référence
reproduite par la référence Python et, ligne à ligne, par la session C++~; chaque message de cette trace analysé et reconstruit
octet pour octet en Rust~; \omterm{def:ll:protocols-i-fix:session}{heartbeat}, demande de test et délai dépassé~; déconnexion sur un numéro de séquence trop bas~; mode de réinitialisation.
\bfield{Pour aller plus loin} Un magasin de messages persistant (le journal binaire du chapitre~23), un Reject pour les messages mal formés, des sessions FIXT~1.1
transportant des messages applicatifs FIX~5.0, et la couche session en Rust.
\end{build}

\begin{omsources}
\item FIX Trading Community, dictionnaire des messages et des champs FIX Latest (Extension Pack 312), tel que publié par le
navigateur Orchimate (champs 7, 8, 9, 10, 16, 34, 35, 36, 43, 49, 52, 56, 108, 112, 123~; messages 0, 1, 2, 4, 5, 8, A, D, 9).
\item FIX Trading Community, spécification de la \emph{FIX Performance Session Layer (FIXP)} (GitHub, versions 1.0 et 1.1).
\item Wikipédia, «~Financial Information eXchange~» (historique~; l'exemple travaillé de longueur du corps et de somme de contrôle).
\item Coinbase Developer Documentation, Exchange FIX API~: connectivity.
\end{omsources}

\section{Exercices}

\begin{exercise}[$\star$]\label{exo:ll:protocols-i-fix:1}
Calculer la longueur du corps et la somme de contrôle du \omterm{def:ll:protocols-i-fix:session}{heartbeat} \texttt{8=FIX.4.4|9=?|35=0|49=A|56=B|34=2|52=20260925-14:30:00.000|10=?|},
où \texttt{|} est le SOH, puis vérifier le résultat avec \texttt{firm\_fixengine.encode}.
\end{exercise}

\begin{exercise}[$\star$]\label{exo:ll:protocols-i-fix:2}
Avec un intervalle de \omterm{def:ll:protocols-i-fix:session}{heartbeat} de 30 secondes et les règles de la construction, le courtier se tait à 10:00:00. Quand la session envoie-t-elle une
demande de test, et quand déclare-t-elle la connexion morte si rien ne revient~?
\end{exercise}

\begin{exercise}[$\star$]\label{exo:ll:protocols-i-fix:3}
La session attend le numéro de séquence 52 et reçoit 57. Qu'envoie-t-elle, que fait-elle du message 57, et que doit
renvoyer le courtier~?
\end{exercise}

\begin{exercise}[$\star\star$]\label{exo:ll:protocols-i-fix:4}
Donner deux messages différents, bien formés, de même longueur et de même somme de contrôle, et expliquer à l'aide de la
\cref{prop:ll:protocols-i-fix:checksum} pourquoi le protocole accepte ce risque.
\end{exercise}

\begin{exercise}[$\star\star$]\label{exo:ll:protocols-i-fix:5}
Pourquoi les \omterm{def:ll:protocols-i-fix:session}{heartbeats} et les demandes de test sont-ils remplacés par un \omterm{def:ll:protocols-i-fix:session}{comblement de trou} au lieu d'être renvoyés~? Que se passerait-il s'ils étaient renvoyés~?
\end{exercise}

\begin{exercise}[$\star\star$]\label{exo:ll:protocols-i-fix:6}
Une passerelle reçoit 200\,000 rapports d'exécution par seconde en pointe. À partir des coûts mesurés, quelle fraction d'un cœur représente
l'analyse avec la vue sans copie, avec l'analyseur fondé sur une table et avec la référence Python~?
\end{exercise}

\begin{exercise}[$\star\star\star$]\label{exo:ll:protocols-i-fix:7}
\emph{Programmation.} Faire en sorte que les sessions Python et C++ répondent à un message dont l'étiquette n'est pas un nombre par un Reject (35=3) portant
RefSeqNum (45), sans avancer au-delà de ce message, puis ajouter ce cas à la conversation de référence pour que les deux traces restent concordantes.
\end{exercise}

\begin{exercise}[$\star\star\star$]\label{exo:ll:protocols-i-fix:8}
\emph{Trouver la faille.} «~Pour garder les choses simples, notre moteur remet à 1 les deux numéros de séquence à chaque reconnexion, donc il n'a jamais
besoin de logique de renvoi.~»
\end{exercise}

\section{Problème~: le trou du week-end}

\begin{problem}[{Problème du week-end --- resynchronisation après une déconnexion}]\label{pb:ll:protocols-i-fix:1}
Le moteur d'un courtier conserve dans son magasin les messages d'une session dont la connexion avec nous est coupée, et les renvoie quand
nous nous reconnectons. Supposer des messages de 205 octets, une liaison de 1~Gb/s avec un aller-retour de \qty{1}{\milli\second}, un intervalle de \omterm{def:ll:protocols-i-fix:session}{heartbeat} de
30 secondes, et les coûts d'analyse mesurés (\texttt{measured\_fix.csv}). Considérer deux déconnexions~: une déconnexion chargée, une liaison qui
meurt silencieusement pendant que le courtier envoie 200 rapports d'exécution par seconde, et une déconnexion calme, 48 heures pendant le week-end avec un
message applicatif par heure et un \omterm{def:ll:protocols-i-fix:session}{heartbeat} chaque fois que 30 secondes s'écoulent sans message.

\medskip\noindent\textbf{Partie I --- La panne chargée.}
\begin{enumerate}
\item Combien de temps la session de la construction met-elle à remarquer la liaison morte, et pourquoi~?
\item Combien de messages le courtier doit-il renvoyer (\texttt{ll\_fix.sender\_stream})~? Combien sont des \omterm{def:ll:protocols-i-fix:session}{heartbeats}~?
\item Combien d'octets cela représente-t-il, et combien de temps prennent-ils sur le fil~?
\item Combien de temps la resynchronisation prend-elle avec le moteur C++, et avec la référence Python~?
\end{enumerate}

\medskip\noindent\textbf{Partie II --- Le week-end calme.}
\begin{enumerate}[resume]
\item Combien de messages attendent-ils pour nous dans le magasin du courtier au bout de 48 heures~? Combien sont des messages applicatifs~?
\item Combien de messages contient la réponse à notre \omterm{def:ll:protocols-i-fix:session}{demande de renvoi}, avec \omterm{def:ll:protocols-i-fix:session}{comblement de trou} et sans~?
\item Combien de temps la resynchronisation prend-elle avec la référence Python, dans chaque sens~?
\item Pourquoi le \omterm{def:ll:protocols-i-fix:session}{comblement de trou} ne change-t-il rien dans le cas chargé~?
\end{enumerate}

\medskip\noindent\textbf{Partie III --- La conception.}
\begin{enumerate}[resume]
\item Pourquoi les messages d'administration ne peuvent-ils jamais être renvoyés tels quels~?
\item Que doit faire l'application d'un rapport d'exécution renvoyé et marqué PossDupFlag~?
\item Une plateforme qui ne prend en charge aucune \omterm{def:ll:protocols-i-fix:session}{demande de renvoi} ouvre une nouvelle session à chaque connexion. Comment apprend-on les exécutions qui
se sont produites pendant qu'on était déconnecté~?
\item Que doit persister le moteur pour survivre à son propre redémarrage sans trou~?
\end{enumerate}

\medskip\noindent\textbf{Partie IV --- Le verdict.}
\begin{enumerate}[resume]
\item Énoncer le \emph{résultat nommé}~: messages renvoyés et temps de resynchronisation dans les deux cas, avec et sans \omterm{def:ll:protocols-i-fix:session}{comblement de trou}.
\item Un intervalle de \omterm{def:ll:protocols-i-fix:session}{heartbeat} de 1 seconde aiderait-il le cas chargé~? À quel coût~?
\item Et si le magasin du courtier ne gardait que ses 10\,000 derniers messages~?
\item La stratégie doit-elle agir sur les rapports d'exécution renvoyés comme s'ils venaient de se produire~?
\item Comment tester la logique de renvoi avant de se connecter à un courtier~?
\item La référence Python peut-elle suivre la session chargée en fonctionnement normal~?
\item Pour lequel des deux cas faut-il concevoir le moteur~?
\item En une phrase~: à quoi sert la couche session~?
\end{enumerate}
\end{problem}

\section{Questions d'entretien}

\begin{interviewq}[$\star$ \iqroles{développeur}]\label{iq:ll:protocols-i-fix:1}
Décrivez la structure d'un message FIX. Comment un récepteur sait-il où un message se termine~?
\end{interviewq}

\begin{interviewq}[$\star\star$ \iqroles{développeur}]\label{iq:ll:protocols-i-fix:2}
Décrivez ce qui se passe lorsqu'une \omterm{def:ll:protocols-i-fix:session}{session FIX} détecte un trou dans les numéros de séquence entrants.
\end{interviewq}

\begin{interviewq}[$\star\star$ \iqroles{développeur}]\label{iq:ll:protocols-i-fix:3}
Comment écririez-vous un analyseur FIX qui n'alloue pas~? Que renvoie-t-il~?
\end{interviewq}

\begin{interviewq}[$\star\star$ \iqroles{développeur}]\label{iq:ll:protocols-i-fix:4}
À quoi servent les \omterm{def:ll:protocols-i-fix:session}{heartbeats} et les demandes de test, et comment choisiriez-vous l'intervalle de \omterm{def:ll:protocols-i-fix:session}{heartbeat}~?
\end{interviewq}

\begin{interviewq}[$\star\star$ \iqroles{développeur}]\label{iq:ll:protocols-i-fix:5}
Pourquoi beaucoup de plateformes utilisent-elles des protocoles binaires pour la saisie d'ordres plutôt que le tag-value de FIX~? Que sacrifient-elles~?
\end{interviewq}

\begin{interviewq}[$\star\star\star$ \iqroles{développeur}]\label{iq:ll:protocols-i-fix:6}
Votre moteur a redémarré en pleine journée et le courtier vous dit que votre numéro de séquence est trop bas. Que s'est-il passé, et comment
reprenez-vous sans perdre ni dupliquer d'exécutions~?
\end{interviewq}
